%% kapa: exam=1 round=24 year=2013 period=2 form=A subject=appraisal_law subject_ko=부동산관계법규 seq=8 number=48
%% source: Q-Net 감정평가사 기출문제 (cst00309) artlSeq=5073802 file=extracted/1008/5073802/2013년 제24회 감정평가사 1차시험 기출문제/2교시(A형).hwp
%% conversion: kordoc 4.24.51 -> tools/gen1.py (auto); review=proofread-sonnet; flags=none
\begin{question}{48}
국토의 계획 및 이용에 관한 법령상 도시ㆍ군계획시설부지의 매수청구에 관한 설명으로 옳은 것은?
\begin{choices}
  \item 매수청구대상에는 도시ㆍ군계획시설의 부지로 되어 있는 토지 중 지목이 잡종지인 토지도 포함된다.
  \item 시장 또는 군수가 해당 도시ㆍ군계획시설사업의 시행자로 정하여진 경우에는 시장 또는 군수가 매수의무자이다.
  \item 매수의무자가 지방공사인 경우 매수대금은 채권으로 지급할 수 있다.
  \item 매수의무자는 매수청구대상 토지에 대한 매수청구를 받은 날부터 3개월 이내에 매수 여부를 결정하여 토지소유자에게 알려야 한다.
  \item 매수의무자가 매수하지 아니하기로 결정한 경우 매수청구를 한 토지소유자는 개발행위허가 없이 그 토지에 공작물을 설치할 수 있다.
\end{choices}
\end{question}
\answer{2}
